% Generated by tools/edn_latex.py from reports/edn.md.
% Do not edit: change reports/edn.md and run `make edn`.
\ednentry{
  id = {EDN-65},
  title = {The Pynblint policy: enforce what can fire on a committed notebook, and show execution by running it},
  date = {2026-10-01},
  milestone = {M3: Quality Assurance},
  topic = {Static Analysis, Testing Strategy},
  participants = {@lukas2510, decided by the repository owner},
  decision = {Of Pynblint 0.1.6's 22 rules, 13 are enforced at their default thresholds and 9 are excluded, each with its reason, in the table in \href{https://github.com/mlops-2627q1-mds-upc/MLOPS_Recommenditos/blob/main/docs/docs/notebooks.md\#pynblint-rules}{Notebooks}.
That table is the configuration: the gate reads it, and refuses to run unless it classifies exactly the rules the installed Pynblint has.
The three rules that read execution counts are excluded, because nbstripout removes those counts on every commit, and that a notebook runs top to bottom is shown by running it with \texttt{make notebook-\allowbreak{}run} instead.
The five repository-level rules are not run, each with the reason and with what enforces the same property here.
Pynblint lints every notebook git would commit, one file at a time, and two rules of our own require every notebook to sit directly in \texttt{notebooks/\allowbreak{}} under the Cookiecutter name \texttt{<number>.\allowbreak{}<version>-\allowbreak{}<initials>-\allowbreak{}<description>.\allowbreak{}ipynb}, and to carry none of the flags that make nbstripout keep its outputs.
There is no warning tier and no per-notebook waiver: a finding fails the build, and a rule that is wrong for us is excluded in the table, with its reason, in a reviewed pull request.},
  rationale = {\emph{Alternatives considered:}
\begin{itemize}[leftmargin=*, topsep=2pt]
\item \textbf{Option A (chosen): exclude the execution rules, keep nbstripout as it is, and evidence execution by running the notebook.}
\newline Pros: nbstripout 0.9.1 sets every \texttt{execution\_\allowbreak{}count} to null, so on a committed notebook \texttt{non-\allowbreak{}executed-\allowbreak{}notebook} always fires, \texttt{non-\allowbreak{}linear-\allowbreak{}execution} can never fire, and \texttt{non-\allowbreak{}executed-\allowbreak{}cells} fires only on an empty code cell and then flags every other code cell; enforcing them would either keep the build red forever or claim checks that do not happen.
Stripping also keeps rows of a dataset with personal data out of git, which is more than diff hygiene.
The profiling notebook's summary cell fails when a dataset card figure does not reproduce, so the drill checks the card as well as the code.
\newline Cons: execution is not gated in CI; \texttt{make notebook-\allowbreak{}run} is a local drill, because the notebook needs the 548 MB source file.
\item \textbf{Option B: \texttt{nbstripout -\allowbreak{}-\allowbreak{}keep-\allowbreak{}count}, so the three rules have something to read.}
\newline Pros: three more live rules.
\newline Cons: execution counts churn in every diff, and they record what one kernel once did, not that a fresh kernel runs the notebook: an edited cell that was never re-run keeps its old count.
Every edit would need a full re-run on the source file before committing, and an editor that writes no counts trips the rules.
\item \textbf{Option C: scan the repository root and run the repository rules.}
\newline Pros: the literal reading of ``notebook and repository QA''.
\newline Cons: \texttt{repository-\allowbreak{}not-\allowbreak{}versioned} looks for a \texttt{.\allowbreak{}git} directory and fires in every git worktree, which this project uses for parallel work; locally the scan walks \texttt{.\allowbreak{}venv} and other worktrees and lints other branches' stale notebook copies; \texttt{test-\allowbreak{}coverage-\allowbreak{}data-\allowbreak{}not-\allowbreak{}available} depends on whether the tests ran in that checkout.
Each repository rule is about a property something else here already enforces more strongly, which the table says rule by rule.
\item \textbf{Option D: scan the \texttt{notebooks/\allowbreak{}} directory.}
\newline Pros: one command, no file list.
\newline Cons: three structural false positives to exclude, untracked scratch files get linted, and a notebook committed anywhere else escapes the scan.
\item \textbf{Option E: the policy in a commented \texttt{.\allowbreak{}pynblint} file.}
\newline Pros: Pynblint's own configuration format.
\newline Cons: it is read only from the current directory and is overridden by environment variables, its reasons are comments nothing checks, and a misspelt rule name is accepted silently, so an exclusion with a typo excludes nothing and says nothing.
\item \textbf{Option F: raise \texttt{cell-\allowbreak{}too-\allowbreak{}long} or \texttt{notebook-\allowbreak{}too-\allowbreak{}long} so the first notebook passes.}
\newline Pros: no restructuring.
\newline Cons: it is the threshold adjusted to its result that EDN-62 rejects; the reason would be the notebook, not the rule.
\item \textbf{Option G: report the excluded rules as non-blocking warnings.}
\newline Pros: nothing is hidden.
\newline Cons: Pynblint has no severities, and a warning tier is exactly the ``findings nobody acts on'' the issue warns about.
\item \textbf{Option H: execute the notebook in CI, with the source file cached under its MD5.}
\newline Pros: the execution evidence that options A and B argue about would be automated, and because the summary cell fails on a mismatch, every pull request would check the dataset card against the data.
\newline Cons: a cold cache downloads 548 MB from Zenodo, so a third party's server would decide whether our CI is green, which is a flaky test by construction; the job needs the full project environment, about a minute and 2.8 GB of RAM; and GitHub evicts a cache unused for seven days, so the download recurs.
It stays open: if the raw data ever comes from our own DVC remote instead, most of the cons go away.
\end{itemize}
\par
\emph{Why:} A rule is enforced where it can fire on what we commit, and excluded where it cannot, with the reason written next to it.
That principle is what keeps the gate honest in both directions: nothing it reports is noise, and nothing it lists as checked is unchecked.
It also bit on the first run: the dataset card's profiling notebook, as first drafted, had five code cells over 30 lines, and they were restructured into shorter cells with precomputed values rather than the threshold being raised, so the notebook stayed within the 50-cell limit as well.
On the stripped copy that is committed, the only rule that fired was \texttt{non-\allowbreak{}executed-\allowbreak{}notebook}, which is the measurement behind excluding the execution rules.
Making the documentation table the configuration follows the pattern of the requirement matrix, which reads the requirements and the specification rather than a copy of them: the reasons cannot drift from what the gate does, and a rule a future Pynblint adds or renames stops the build until someone decides about it.
The table names every rule, enforced or excluded, and the gate refuses one that does not classify exactly the installed rules, enforces nothing, or enforces a repository rule that would never run, so a misspelt slug, a rule a new release adds or an empty policy stops the gate rather than silently changing what is checked.},
  ai = {Information seeking, Alternative generation, Alternative assessment, Recommendation, Solution generation},
  response = {Accepted},
  assessment = {AI read the Pynblint 0.1.6 source and tested every rule on purpose-built notebooks, stripped and unstripped, which is how the interaction with nbstripout was found rule by rule rather than assumed; it also found the blind spots the table now states as caveats.
It proposed the policy, the docs-as-configuration design and the location rule, and laid out options A to G; the owner chose A for the execution rules and the docs table as the policy.
The notebook itself was drafted by AI, and its first Pynblint run found five code cells over 30 lines; they were restructured rather than the threshold raised, so the policy applied to AI-written code as to any other.
An AI review of the documents then found two holes in the first version of this policy, both closed before the pull request was opened: nothing stopped a \texttt{keep\_\allowbreak{}output} flag from carrying outputs past nbstripout, and \texttt{make notebook-\allowbreak{}run}, the stand-in for the excluded execution rules, passed even when a dataset card figure did not reproduce.},
  aievidence = {Claude Code session on 2026-10-01 implementing issue \#41: a rule-by-rule experiment on stripped and unstripped sample notebooks, an architecture pass that produced the policy table and options A to G, the owner's choice among them, and a document review that added option H and the two fixes above.},
  evidence = {\href{https://github.com/mlops-2627q1-mds-upc/MLOPS_Recommenditos/blob/main/docs/docs/notebooks.md}{Notebooks}, the rule table; \href{https://github.com/mlops-2627q1-mds-upc/MLOPS_Recommenditos/blob/main/notebooks/1.0-lh-dataset-card-profiling.ipynb}{\texttt{notebooks/\allowbreak{}1.\allowbreak{}0-\allowbreak{}lh-\allowbreak{}dataset-\allowbreak{}card-\allowbreak{}profiling.\allowbreak{}ipynb}}; \texttt{tests/\allowbreak{}test\_\allowbreak{}notebook\_\allowbreak{}lint.\allowbreak{}py}; \href{https://github.com/mlops-2627q1-mds-upc/MLOPS_Recommenditos/blob/main/docs/docs/specification.md}{specification} NFR-07; \hyperref[EDN-64]{EDN-64}; \href{https://github.com/mlops-2627q1-mds-upc/MLOPS_Recommenditos/issues/41}{issue \#41}, \href{https://github.com/mlops-2627q1-mds-upc/MLOPS_Recommenditos/pull/67}{PR \#67}.},
}
